\lecture{36}{Wed 07 Dec 2022 11:33}{}

\begin{theorem}(Uniform convergence and integrability)
	Suppose $f_n: [a,b] \to \mathbb{R}$ are R-S integrable wrt $g$, $g$ monotone increase on $[a,b]$. If $f(x) = \sum_{n=1}^{\infty} f_n(x)$ is uniformly convergent on $[a,b]$ then $f$ is R-S integrable wrt $g$ on $[a,b]$ and 
	\[
		\int _a^b f dg = \sum_{n=1}^{\infty} \int _a^b f_n dg
	.\] 
\end{theorem}

\begin{theorem}(Uniform convergence and Differentiability)
	Suppose $f_n: [a,b] \to \mathbb{R}$ (continuously) differentiable on $[a,b]$ such that 
	\begin{itemize}
		\item $\sum_{n=1}^{\infty} f'_n(x)$ is uniformly cnvergent on $[a,b]$
		\item  $\sum_{n=1}^{\infty} f_n(x_0)$ is convergent for some $x_0 \in [a,b]$.
	\end{itemize}
	Then, $f(x) = \sum_{n=1}^{\infty} f_n(x)$ is convergent uniformly on $[a,b]$, and $f'(x) = \sum_{n=1}^{\infty} f'_n(x)$, $x \in [a,b]$.
\end{theorem}
This justifies term-by-term differentiation.

\begin{remark}
	This theorem is true without assuming continuously differentiable. However the proof we give will depend on this assumption.
\end{remark}

\begin{proof}
	From assumptions and Newton Leibniz formula, we have
	\[
		f_n(x) = f_n(x_0) + \int _{x_0}^x f'_n(x) dx
	.\] 
	Taking sums,
	\[
		\sum_{n=1}^{\infty} f_n(x) = \sum_{n=1}^{\infty} f_n(x_0) + \sum_{n=1}^{\infty} \int _{x_0}^x f'_n(t) dt
	.\] 
	Applying definition of $f(x)$ and uniform convergence, we have
	\[
		f(x) = \sum_{n=1}^{\infty} f_n(x_0) + \int _{x_0}^t \sum_{n=1}^{\infty}f'_n(t)dt 
		=:  f(x_0) + \int _{x_0}^x g(t) dt
	.\] 
	By the differentiation theorem,
	\[
		f'(x) = g(x) \qquad 
		\text{ for any }  x \in [a,b]
	.\] 
\end{proof}

\begin{example}
	Show that
	\[
		f(x) = \sum_{n=1}^{\infty} \frac{(nx) }{n^2}
		\qquad x \in \mathbb{R}
	\] 
	where
	\[
		(\alpha) = \alpha - [\alpha] \in [0,1)
	\]
	is the fractional part of $\alpha$,
	is discontinuous at all rational points, but is Riemann integrable on any finite interval.
\end{example}
	By Weierstrass M-test, $\sum f_n(x)$ is uniformly convergent. 

	$f_n$ has finitely many discontinuity, so $f_n$ is Riemann integrable, hence $f$ is Riemann integrable.

	$f_n$ not continuous at all $c$, but the ine sided limits exist at every point. Define
	\[
		j_n(c) = f_n(c+) - f_n(c-)
	.\] 
	be the \textit{jump} at $c$. 

	By uniform continuity,
	\[
		f(c+) = \lim_{x \to c+} \sum f_n(x) = \sum \lim_{x \to c+} f_n(x)
	.\] 

	Similarly we define $f(c-)$. Now we have
	\[
		f(c+) - f(c-) = \sum_{n=1}^{\infty} j_n(c) < 0 \iff  c \in \mathbb{Q}
	.\] 

\subsection{Power Series}

A power series is a series of the form
\[
	f(x) = \sum_{n=0}^{\infty} a_n x^n
.\] 

\begin{theorem}
	Let $\rho = \limsup_{n \to \infty} |a_n|^{\frac{1}{n}} \in [0, \infty ]$ and $R = \frac{1}{\rho}$ in the extended sense. 

	Then
	\begin{enumerate}
		\item for any $|x| < R$ the series is convergent.
		\item for any $|x| \le  C < R$ the series is uniformly convergent.
		\item If $|x| >R$ the series is divergent.
		\item for $|x| = R$ the series may or may not diverge.
	\end{enumerate}
	$R$ is called \logseq{Radius of Convergence}.
\end{theorem}
\begin{proof}
	Apply Root test to $a_n x^n$. Thus
	\[
		\limsup |a_n x^n|^{\frac{1}{n}} = |x| \limsup |a_n|^{\frac{1}{n}} =  \frac{|x|}{R}< 1 \iff |x| > R
	.\]

	To prove the uniformly convergent part, let $x \in [-r, r] $ and
	\[
		|a_n x^n| \le |a_n| r^n
	.\] 
	and
	\[
		\sum |a_n| r^n < \infty  \qquad \text{ by Root Test}
	.\]
	Hence $\sum a_n x^n$ uniformly converge by Weierstrass M-test.
\end{proof}

\begin{theorem}
	If $f$ is a power series as above, and $R$ is as above, then $f$ is infinitely differentiable in $(-R, R)$ and the series can be differentiated term-by-term.
\end{theorem}
\begin{proof}
	Consider the series
	$g(x) = \sum_{k=1}^{\infty} k a_k x^{k-1}$. $g(x)$ converges iff $x g(x) = \sum_{k=0}^{\infty} k a_k x^k $ converges. We claim this sum has the same radius of convergence as $g$. This is true since $k^{\frac{1}{k}}\to  1$. 

	Besides, we know $f$ is convergent at $0$. Apply differentiation theorem,
	$f'(x) = \sum_{k=0}^{\infty} k a_k x^{k-1}$.
\end{proof}
\begin{remark}
	From term-of-term differentiation we have
	\[
		a_n = \frac{f^{(n)}(0)}{n!}
	.\] 
	and $f(x)$ is a \logseq{Taylor series}.
\end{remark}




